\scanpage{82}
\mspage{(40)}
\[
\tag{105}
\diagram{\begin{tikzcd}[ampersand replacement=\&, column sep=small,row sep=large]
\Pi M^*_{G,\deleted{\unsure{G'}},G''}
 \arrow[rr,dashed] \arrow[d,dashed]
 \arrow[dd,bend right=35,"\varphi_{G,G''}"']
 \arrow[rrr,bend left=22,"\gamma_{G,G''}"]
\&\& \Pi M^*_{G',G''} \arrow[r,"\gamma_{G',G''}"] \arrow[d] \& \Pi M_{G''} \\
\Pi M^*_{\underline G,\underline G'} \arrow[r,"\gamma_{G,G'}"] \arrow[d]
\& \deleted{\unsure{\Pi M_{G'}}} \arrow[r,"\deleted{\mathrm{can}}"]
\& \begin{gathered}\Pi M_{G'\deleted{,G''}}\\(\simeq\Pi(M_{G'},\Gamma^!_{G',G''}))\end{gathered}
\& \text{(103) pour }(G',G'')\\
\Pi M_{G\deleted{,\unsure{G'}}} \& \text{(103)}
\end{tikzcd}}
\]
\margin{(2 cart)}
\deleted{$(\unsure{\simeq\Pi M})$}

Dans ce diagramme, on a \deleted{\unsure{enlevé|oublié}} laissé tomber les soulignés pour les MD-graphes, \deleted{et on a remplacé la notation $\Pi(M_{G,G'},\Gamma^!_{G,G'})$ par la notation plus légère $\Pi M_{G,G'}$, de $=$ pour $\Pi M_{G',G''}$.} On \deleted{\unsure{tire de (105)}}\deleted{a défini} \unsure{retrouve}
\[
\tag{106}
\Pi M^*_{G,G''\deleted{,\unsure{G'}}}
\]
dans (105) comme un produit 2-cartésien --- c'est donc un groupoïde connexe, ext. de $\Pi M^*_{G,G'}$ par un $\mathbf Z^{C'}$ (tordu), et il se trouve que c'est aussi, \unsure{en tant qu'ext.} de $\Pi M_{G\deleted{,\unsure{G'}}}$ (de noyau une ext. de $\mathbf Z^C$ tordu par $\mathbf Z^{C'}$ tordu) une extension de $\Pi M_{G\deleted{,\unsure{G'}}}$ par un $\mathbf Z^{C\amalg C'}$ tordu.

De façon plus précise, on tire de (105) le diagramme
\[
\tag{107}
\diagram{\begin{tikzcd}[ampersand replacement=\&, column sep=large,row sep=large]
\Pi M^*_{G,\deleted{\unsure{G'}},G''} \arrow[r,dash,"\gamma_{G,\deleted{\unsure{G'}},G''}"] \arrow[d,dash,"\varphi_{G,G''\deleted{,\unsure{G'}}}"] \& \Pi M_{G''} \\
\Pi M_{G\deleted{,\unsure{G'}}}
\end{tikzcd}}
\]
\margin{81}
